\documentclass[12pt]{article}
\usepackage{graphicx} %package to manage images
\graphicspath{ {./figures/} }
\usepackage{caption}
\usepackage[font=scriptsize]{subcaption}
\captionsetup[figure]{labelsep=none}
\captionsetup[table]{labelsep=none}
\usepackage{bbm}
\usepackage{amsmath}
\usepackage{import}
\usepackage{array}
\usepackage{booktabs}     % no extra vertical space between paragraphs
\usepackage{indentfirst}
\usepackage{afterpage}
\usepackage{floatrow}
\usepackage{pdflscape}
\usepackage{soul}
\usepackage{float}
\usepackage{adjustbox}
\usepackage{longtable}
\usepackage{caption}
\usepackage{setspace}
\usepackage{afterpage}
\usepackage[margin=1in]{geometry}
\usepackage[round]{natbib}
\usepackage{hyperref}
\usepackage{titlesec}
\usepackage{threeparttable}
\usepackage{setspace}
\usepackage{float}
\floatstyle{plaintop}
\restylefloat{table}
\usepackage{rotating}
\usepackage{ragged2e}
\usepackage{authblk}

% Optional styling for nicer author/affil typography
\renewcommand\Authfont{\normalsize\bfseries}
\renewcommand\Affilfont{\small\normalfont}
\setlength{\affilsep}{0.3em} % space between author and affiliation

% \author[1,2,3]{Diane Coffey\thanks{Corresponding author: coffey@utexas.edu}}

% \author[1]{Siddhant Pandit}

% \affil[1]{Population Research Center, UT Austin}
% \affil[2]{Department of Sociology, UT Austin}
% \affil[3]{r.i.c.e., a research institute for compassionate economics}

% \setlength{\parindent}{0pt}   % standard indent (~0.25 inch)
% \setlength{\parskip}{1em}       % no extra vertical space between paragraphs

\titleformat{\section}
  {\normalfont\bfseries\fontsize{15.5}{16}\selectfont}  % Bold 14pt
  {}{0pt}{}  % No numbering

\titleformat{\subsection}
  {\normalfont\bfseries\fontsize{12}{12}\selectfont}  % Bold 12pt
  {}{0pt}{}  % No numbering

% \titleformat{\subsubsection}
%   {\normalfont\itshape\fontsize{12}{14}\selectfont}  % Italic 12pt
%   {}{0pt}{}  % No numbering
  
\begin{document}
% \RaggedRight
 
\section*{Disparities in Maternal Nutrition among Indian Social Groups: The Roles of Fertility Patterns and Household Wealth} 

\subsection*{Relevance to Demography Statement}

\vspace{0.25cm}
% \newline

\noindent This study uses decomposition methods and the latest Demographic and Health Survey to describe variation in maternal nutrition in India, where one in five women begins pregnancy underweight. We contribute to the health disparities literature the first estimates of pregnancy underweight by caste and religion.  We ask whether differences in fertility patterns explain why women from oppressed castes are more likely to be underweight before pregnancy.  We find that parity progression and birth spacing explain only a small fraction of the gaps we observe.  Rather, social group differences in prepregnancy underweight are substantially explained by differences in household wealth. 

\newpage
\doublespacing

\end{document}